% Part: incompleteness
% Chapter: arithmetization-syntax
% Section: coding-formulas

\documentclass[../../../include/open-logic-section]{subfiles}

\begin{document}

\olfileid{inc}{art}{frm}
\olsection{Pangodean \printtoken{P}{formula}}

Sawisé kita netepake relasi $\fn{Term}(x)$
kanthi rekursif primitif, relasi iki bisa
digunakake kanggo netepake relasi kang
cocog kanggo !!{formula}s, yaiku
$\fn{Frm}(x)$, kanthi rekursif primitif.

\begin{prop}
Relasi $\fn{Atom}(x)$, kang bener yen lan
mung yen $x$ iku bilangan G\"odel saka
!!{formula} atomik, iku rekursif primitif.
\end{prop}

\begin{proof}
Bilangan $x$ iku bilangan G\"odel saka
!!{formula} atomik yen lan mung yen
salah siji saka kahanan iki bener:
\begin{enumerate}
\item Ana $n$, $j < x$, lan $z < x$
  saengga kanggo saben $i < n$,
  $\fn{Term}((z)_i)$ lan $x = $
\[
\Gn{\Obj P^n_j(} \concat \fn{flatten}(z) \concat \Gn{)}.
\]
\item Ana $z_1, z_2 < x$ saengga
  $\fn{Term}(z_1)$, $\fn{Term}(z_2)$,
  lan $x = $
\[
\Gn{{\eq}(} \concat z_1 \concat \Gn{,} \concat z_2 \concat{\Gn{)}}.
\]
  \tagitem{prvFalse}{$x = \Gn{\lfalse}$.}{}
  \tagitem{prvTrue}{$x = \Gn{\ltrue}$.}{}
\end{enumerate}
\end{proof}

\begin{prop}
\ollabel{prop:frm-primrec}
Relasi $\fn{Frm}(x)$, kang bener yen lan
mung yen $x$ iku bilangan G\"odel saka
!!a{formula}, iku rekursif primitif.
\end{prop}

\begin{proof}
Urutan simbol~$s$ iku !!a{formula} yen lan
mung yen ana urutan pambentukan~$s_0$,
\dots, $s_{k-1} = s$ saka !!{formula}
kang nyathet carane~$s$ diwangun saka
!!{formula}s atomik miturut aturan
pambentukan. Kode saben $s_i$
(lan uga kode saka urutan
$\tuple{s_0, \dots, s_{k-1}}$)
luwih cilik tinimbang kode~$x$ saka~$s$.
% OLPL-169: Sumber ngaku kode urutan pambentukan luwih cilik tinimbang x; kanggo formula atomik kanthi urutan siji [s], kodene tuple{x}=2^(x+1)>x. Klaim iki perlu didandani ing sumber.
\end{proof}

\begin{prob}
Wenehana bukti rinci kanggo
\olref[inc][art][frm]{prop:frm-primrec}
kanthi ngetutake pola bukti kapisan saka
\olref[inc][art][trm]{prop:term-primrec}.
\end{prob}

\begin{prop}
\ollabel{prop:freeocc-primrec}
Relasi $\fn{FreeOcc}(x, z, i)$, kang bener
yen lan mung yen simbol kapangka-$i$
ing formula kanthi bilangan G\"odel~$x$
iku kemunculan bebas saka variabel kanthi
bilangan G\"odel~$z$, iku rekursif primitif.
\end{prop}

\begin{proof}
Latihan.
\end{proof}

\begin{prob}
Buktèkna \olref[inc][art][frm]{prop:freeocc-primrec}.
Kowé bisa nggunakake fakta yen saben
subuntai saka !!a{formula} kang uga
!!a{formula} iku sub-!!{formula} saka
formula mau.
\end{prob}

\begin{prop}
Sifat $\fn{Sent}(x)$, kang bener yen lan
mung yen $x$~iku bilangan G\"odel saka
!!a{sentence}, iku rekursif primitif.
\end{prop}

\begin{proof}
!!{sentence} iku !!a{formula} tanpa
kemunculan bebas !!{variable}s.
Mula $\fn{Sent}(x)$ bener yen lan mung yen
\begin{multline*}
\bforall{i<\len{x}}{\bforall{z<x}{}}\\
(\bexists{j<z}{z=\Gn{\Obj v_j}} \lif \lnot\fn{FreeOcc}(x,z,i)).
\end{multline*}
\end{proof}


\end{document}
